\documentclass[11pt,a4paper]{article}
\usepackage[utf8]{inputenc}
\usepackage[T1]{fontenc}
\usepackage[ngerman]{babel}
\usepackage{geometry}
\usepackage{microtype}
\usepackage{amsmath,amssymb,mathtools}
\usepackage{booktabs,longtable,array}
\usepackage{enumitem}
\usepackage{xcolor}
\usepackage{hyperref}
\usepackage{listings}
\usepackage{fancyhdr}
\usepackage{titlesec}
\usepackage{tikz}
\usetikzlibrary{arrows.meta,positioning,shapes.geometric,fit}

\geometry{margin=2.2cm}
\hypersetup{colorlinks=true, linkcolor=black, urlcolor=blue, citecolor=black}
\pagestyle{fancy}
\fancyhf{}
\lhead{PSE-TRAVERSE-PROJECTION-02}
\rhead{Hypertorus-Gated Finalization Layer v0.2}
\cfoot{\thepage}
\setlist[itemize]{topsep=3pt,itemsep=2pt,parsep=0pt,leftmargin=*}
\setlist[enumerate]{topsep=3pt,itemsep=2pt,parsep=0pt,leftmargin=*}
\titleformat{\section}{\Large\bfseries}{\thesection}{0.8em}{}
\titleformat{\subsection}{\large\bfseries}{\thesubsection}{0.8em}{}
\titleformat{\subsubsection}{\normalsize\bfseries}{\thesubsubsection}{0.8em}{}

\definecolor{codebg}{RGB}{246,246,246}
\definecolor{rulegray}{RGB}{110,110,110}
\lstdefinelanguage{RustLike}{
  morekeywords={pub,struct,enum,fn,let,mut,impl,trait,where,Option,Vec,BTreeMap,BTreeSet,String,bool,u64,u32,u8,i128,Self,return,match,if,else,true,false},
  sensitive=true,
  morecomment=[l]{//},
  morestring=[b]"
}
\lstset{
  language=RustLike,
  backgroundcolor=\color{codebg},
  basicstyle=\ttfamily\small,
  frame=single,
  rulecolor=\color{rulegray},
  breaklines=true,
  columns=fullflexible,
  keepspaces=true,
  showstringspaces=false
}

\newcommand{\MUST}{\textbf{MUST}}
\newcommand{\MUSTNOT}{\textbf{MUST NOT}}
\newcommand{\SHOULD}{\textbf{SHOULD}}
\newcommand{\MAY}{\textbf{MAY}}
\newcommand{\code}[1]{\texttt{\detokenize{#1}}}

\title{\textbf{PSE-TRAVERSE-PROJECTION-02}\\[0.35em]
Hypertorus-Gated Finalization Layer\\
Formale Implementierungsspezifikation v0.2}
\author{Sebastian Klemm\\Technische Ausarbeitung fuer Coding-Agent-Implementierung}
\date{Mai 2026}

\begin{document}
\maketitle

\begin{abstract}
Diese Spezifikation definiert die normative Gesamtschicht \emph{Hypertorus-Gated Finalization Layer} fuer \code{pse-traverse}. Version 0.2 integriert die Universal Projection Fabric aus v0.1 mit einer deterministischen Hypertorus-Taktung, phasenfensterbasierter Ereignishorizont-Logik, Exclusion/Expansion/Finalization-Matrix, Action-Vector-Bildung, FinalizedEmission-Artefakten und einer streng replayfaehigen Gate-Pipeline. Version 0.2 ist als vollstaendige Nachfolgespezifikation zu behandeln: Eine Implementierung kann v0.1 konzeptionell verwerfen, sofern alle v0.1-Kernfunktionen durch die hier definierten Typen, Operatoren, Reports, Gates und Tests abgedeckt werden. Die Schicht erzeugt keine \code{SemanticCrystal}-Instanzen. Sie erzeugt ausschliesslich finalisierte, gate-gepruefte, content-addressed Emissions- und Zertifikatsobjekte, die optional an die bestehende PSE-Bridge weitergereicht werden.
\end{abstract}

\tableofcontents
\newpage

\section{Status, Geltungsbereich und normative Einordnung}

\subsection{Status}
Dieses Dokument ist eine normative Implementierungsspezifikation. Die Begriffe \MUST, \MUSTNOT, \SHOULD{} und \MAY{} sind verbindlich im Sinne technischer Konformitaet.

\subsection{Versionierung}
\code{PSE-TRAVERSE-PROJECTION-02} ersetzt \code{PSE-TRAVERSE-PROJECTION-01} als normative Zielarchitektur. Version 0.1 bleibt historisch nachvollziehbar, ist fuer Implementierungsentscheidungen jedoch nur noch als abgedeckte Teilmenge zu behandeln.

\begin{center}
\begin{tabular}{p{0.23\linewidth}p{0.68\linewidth}}
\toprule
Version & Rolle \\
\midrule
v0.1 & Universal Projection Fabric: NullCenter, MirrorSeam, Telescope, Loopback, PathExcision, TopologicalWitness, RealityCut, GatedEmission. \\
v0.2 & Vollstaendige Projektions- und Finalisierungsschicht: alle v0.1-Funktionen plus HypertorusScheduler, PhaseHorizonWindow, EpochCounter, ExclusionAxis, ExpansionAxis, FinalizationCore, HolisticMatrixGate und FinalizedEmission. \\
\bottomrule
\end{tabular}
\end{center}

\subsection{Ziel}
Ziel ist eine deterministische, auditierbare, phasengetaktete Projektions- und Finalisierungsschicht fuer stabilisierte Traversal-Zustaende. Ein Zustand darf nur dann externalisiert und finalisiert werden, wenn:

\begin{enumerate}
  \item sein Nullzentrum kanonisch und replayfaehig ist;
  \item seine Projektionen ueber Mirror-Seams zum selben Nullzentrum zurueckfuehren;
  \item ein deklarierter Hypertorus-Phasenhorizont offen ist;
  \item Exclusion, Expansion und Finalization gleichzeitig bestanden sind;
  \item RealityCut, TopologicalWitness, QualityCascade und ReplayReady bestanden sind;
  \item alle gate-relevanten Daten content-addressed und deterministisch serialisiert sind.
\end{enumerate}

\subsection{Nicht-Zweck}
Die Schicht ist nicht:
\begin{itemize}
  \item ein Ersatz fuer \code{pse-core};
  \item ein Erzeuger von \code{SemanticCrystal};
  \item ein LLM-Wrapper;
  \item ein physikalisches Singularitaets-, Kosmologie- oder Raumzeitmodell;
  \item ein Quantum-Hardware-System;
  \item ein Riemann-, Zeta-, Trading-, Hardware-, Fahrzeug- oder Cyber-Defense-Kern;
  \item ein probabilistischer Scheduler ohne Replay-Garantie.
\end{itemize}

\section{Systemrolle innerhalb der Gesamtarchitektur}

\subsection{Schichtmodell}
\begin{longtable}{p{0.28\linewidth}p{0.62\linewidth}}
\toprule
Schicht & Rolle \\
\midrule
\endhead
\code{pse-core} & Evidence-, Gate-, Crystal- und Commit-Kern. Nur diese Schicht darf \code{SemanticCrystal} erzeugen. \\
\code{pse-traverse-core} & \code{ProblemSpec}, \code{FieldCube}, \code{DoFGraph}, \code{CollapsePlan}, \code{Candidate}, \code{GateReport}, \code{CommitOutcome}. \\
\code{pse-traverse-signature} & Operator- und Signaturdiagnostik struktureller Raeume. \\
\code{pse-traverse-dynamics} & Lift/Projection, Field Absorption, Guidance Field, Compressor, TransitionProof. \\
\code{pse-traverse-reactive} & TriTemporalClock, ZeroPointCoupling, ResidualZero, ReactiveGate. \\
\code{pse-traverse-projection-v0.2} & Hypertorus-Gated Finalization Layer: NullCenterProjection, MirrorSeam, HypertorusScheduler, PhaseHorizonWindow, ExclusionAxis, ExpansionAxis, FinalizationCore, HolisticMatrixGate, FinalizedEmission. \\
\bottomrule
\end{longtable}

\subsection{Harte Architekturregel}
\begin{quote}
\textbf{Normative Regel.} Diese Schicht \MUSTNOT{} eine \code{SemanticCrystal}-Instanz erzeugen. Sie \MAY{} ein \code{FinalizedEmission}, ein \code{ProjectionCertificate}, ein \code{FinalizationCertificate} und maschinenlesbare Reports erzeugen. Ein Crystal entsteht ausschliesslich ueber die bestehende PSE-Bridge.
\end{quote}

\section{Reinkristall der Projektion}

\subsection{Neutralisierte Topologie}
Die implementierbare Struktur lautet:
\[
\text{NullCenter}
\rightarrow \text{PhaseCarrierSet}
\rightarrow \text{EventHorizonWindow}
\rightarrow \text{ProjectionCone}
\rightarrow \text{Exclusion/Expansion/Finalization}
\rightarrow \text{FinalizedEmission}.
\]

\begin{center}
\begin{tikzpicture}[node distance=1.4cm,>=Latex, every node/.style={font=\small}]
\node[draw,rounded corners,align=center,fill=gray!8] (n) {NullCenter\\canonical origin};
\node[draw,rounded corners,align=center,fill=gray!8,right=of n] (p) {PhaseCarrierSet\\Hypertorus};
\node[draw,rounded corners,align=center,fill=gray!8,right=of p] (h) {PhaseHorizon\\Gate window};
\node[draw,rounded corners,align=center,fill=gray!8,right=of h] (c) {ProjectionCone\\target views};
\node[draw,rounded corners,align=center,fill=gray!8,below=of c] (m) {HolisticMatrix\\Exclusion + Expansion + Finalization};
\node[draw,rounded corners,align=center,fill=gray!8,left=of m] (e) {FinalizedEmission\\irreversible output node};
\draw[->] (n) -- (p);
\draw[->] (p) -- (h);
\draw[->] (h) -- (c);
\draw[->] (c) -- (m);
\draw[->] (m) -- (e);
\draw[->,dashed] (e) to[bend left=40] node[below]{replay / certificate} (n);
\end{tikzpicture}
\end{center}

\subsection{Singularitaets-View als technisches Modell}
Die bildhafte Formulierung einer First-Person-Sicht aus einer Singularitaet wird rein technisch modelliert:

\begin{longtable}{p{0.32\linewidth}p{0.58\linewidth}}
\toprule
Bildbegriff & Implementierbarer Begriff \\
\midrule
\endhead
Singularitaet & \code{NullCenterObject}: kanonischer Digest-/Self-Metadata-Kern. \\
Ereignishorizont & \code{PhaseHorizonWindow}: phasen- und epochengebundenes Gate-Fenster. \\
Zeitlinien / Phasenlinien & \code{PhaseCarrier}: deterministische Phasentraeger im Hypertorus. \\
Ferne / Ausdehnung & \code{ProjectionCone}: Menge erlaubter Zielraum- und Artefaktprojektionen. \\
Kollaps & \code{FinalizationCore}: irreversible Emissionsentscheidung nach Gate-Pass. \\
\bottomrule
\end{longtable}

Diese Begriffe sind keine Physikbehauptungen. Sie sind Namen fuer Datenstrukturen und Operatoren innerhalb einer deterministischen Kontrollarchitektur.

\section{Mathematische Basisschicht: Phase, Torus und Hypertorus}

\subsection{Zeit zu Phase}
Jeder PhaseCarrier besitzt eine diskrete oder kontinuierlich approximierte Phasenentwicklung. Im Audit-Pfad \MUST{} die diskrete Version verwendet werden.
\[
\theta(k) = \theta_0 + \sum_{j=0}^{k-1} \omega(j) \Delta k.
\]
Bei konstantem Takt:
\[
\theta(k) = \theta_0 + k \omega.
\]
Die normalisierte Phase ist:
\[
\widetilde{\theta}(k) = \theta(k) \bmod 2\pi.
\]

\subsection{Hypertorus}
Ein \(n\)-phasiger Hypertorus ist:
\[
T^n = S^1 \times \cdots \times S^1.
\]
Ein Zustand wird als Phasenvektor repraesentiert:
\[
\Theta(k) = (\theta_1(k),\dots,\theta_n(k)).
\]
Fuer Implementierungen wird die Einbettung in \(\mathbb{R}^{2n}\) verwendet:
\[
U(k)=\left(\cos\theta_1,\sin\theta_1,\dots,\cos\theta_n,\sin\theta_n\right).
\]

\subsection{Epochen und Phasenfenster}
Ein Epochenzaehler ist:
\[
\operatorname{epoch}_i(k)=\left\lfloor\frac{\theta_i(k)}{2\pi}\right\rfloor.
\]
Ein Fensterindikator ist:
\[
\mathbf{1}_{[\alpha,\beta]}(\widetilde{\theta}_i)=
\begin{cases}
1, & \widetilde{\theta}_i \in [\alpha,\beta],\\
0, & \text{sonst.}
\end{cases}
\]

\subsection{Inkommensurable Scheduler}
Fuer mehrere Carrier \(i,j\) \SHOULD{} gelten:
\[
\frac{\omega_i}{\omega_j}\notin \mathbb{Q}
\]
als deklarierte Designintention. Da Irrationalitaet auf Maschinen nur approximiert werden kann, \MUST{} die Implementierung rational kodierte Frequenzen mit deterministischer Perioden- und Kollisionsanalyse verwenden. Der Audit-Pfad \MUST{} die verwendeten rationalen Approximationen speichern.

\section{Normative Begriffe}

\begin{longtable}{p{0.26\linewidth}p{0.64\linewidth}}
\toprule
Begriff & Normative Bedeutung \\
\midrule
\endhead
\code{NullCenterObject} & Kanonischer Ursprung der Projektion, gebunden an Input, RunDescriptor, Operatorversionen und Evidence. \\
\code{PhaseCarrier} & Deterministischer Phasentraeger mit Frequenz, Startphase, Epoche, Quantisierung und Replay-ID. \\
\code{HypertorusScheduler} & Scheduler ueber mehreren PhaseCarriern; erzeugt Fenster, Adressen, Quantisierungsklassen und nicht-wiederholende Zugriffsmuster. \\
\code{PhaseHorizonWindow} & Gate-Fenster, in dem Projektion, Expansion oder Finalisierung erlaubt sein kann. \\
\code{ExclusionAxis} & Gate-Achse zur Eliminierung instabiler Pfade; entspricht neutralisiertem RealityCut/PoR/Telescope. \\
\code{ExpansionAxis} & Gate-Achse zur Erzeugung eines gerichteten Handlungspotentials aus stabiler Resonanz. \\
\code{FinalizationCore} & Entscheidungskern, der gueltige Projektion und Handlungspotential in ein irreversibles Emissionsobjekt ueberfuehrt. \\
\code{HolisticMatrixGate} & Gesamtes Gate: Exclusion, Expansion, Finalization, UPF-Gate, Horizon, Replay und Evidence muessen gleichzeitig bestehen. \\
\code{FinalizedEmission} & Finales, content-addressed, replayfaehiges Emissionsobjekt; kein SemanticCrystal. \\
\bottomrule
\end{longtable}

\section{Kanonisches Datenmodell}

\subsection{Primitive}
\begin{lstlisting}
pub type Hash256 = [u8; 32];
pub type StableId = String;
pub type Fixed = i128; // Q64.64 or project-wide fixed-point type

pub struct EvidenceRef {
    pub kind: String,
    pub address: Hash256,
    pub relation: String,
}
\end{lstlisting}

Alle gate-relevanten numerischen Werte \MUST{} als Fixed-Point, Integer-Rational oder dokumentierte plattforminvariante Repräsentation implementiert werden. Plattform-Floats \MUSTNOT{} Teil des Gate-, Audit- oder Digest-Pfades sein.

\subsection{ProjectionRunDescriptor v0.2}
\begin{lstlisting}
pub struct ProjectionRunDescriptorV2 {
    pub run_id: StableId,
    pub problem_spec_hash: Hash256,
    pub traversal_report_hash: Option<Hash256>,
    pub seed: u64,
    pub operator_versions: BTreeMap<String, String>,
    pub thresholds: ProjectionThresholdsV2,
    pub policies: ProjectionPoliciesV2,
    pub hypertorus: HypertorusSchedulerSpec,
    pub finalization: FinalizationPolicy,
    pub canonicalization_version: String,
}
\end{lstlisting}

\subsection{Thresholds}
\begin{lstlisting}
pub struct ProjectionThresholdsV2 {
    pub max_null_distance: Fixed,
    pub min_mirror_consistency: Fixed,
    pub min_por: Fixed,
    pub min_loopback: Fixed,
    pub min_topological_witness: Fixed,
    pub max_reality_cut_delta: Fixed,
    pub max_path_excision_delta: Fixed,
    pub min_projection_quality: Fixed,
    pub min_exclusion_score: Fixed,
    pub min_expansion_score: Fixed,
    pub min_finalization_score: Fixed,
    pub min_horizon_alignment: Fixed,
    pub max_phase_jitter: Fixed,
    pub max_steps: u64,
}
\end{lstlisting}

\subsection{Policies}
\begin{lstlisting}
pub struct ProjectionPoliciesV2 {
    pub seam_selection: SeamSelectionPolicy,
    pub mount_order: Vec<MountSpec>,
    pub telescope_policy: TelescopePolicy,
    pub condensation_policy: CondensationPolicy,
    pub path_excision_policy: PathExcisionPolicy,
    pub reality_cut_policy: RealityCutPolicy,
    pub emission_policy: EmissionPolicy,
    pub horizon_policy: HorizonPolicy,
    pub finalization_policy: FinalizationPolicy,
}

pub enum SeamSelectionPolicy { Star, Clique, Explicit(Vec<(StableId, StableId)>) }
pub enum EmissionPolicy { HoldOnly, ProjectionBundle, FinalizedEmission }
pub enum HorizonPolicy { StrictWindow, EpochAligned, CarrierQuorum, DisabledForTests }
\end{lstlisting}

\section{Hypertorus Scheduler}

\subsection{Datenstrukturen}
\begin{lstlisting}
pub struct RationalFixed {
    pub numerator: i128,
    pub denominator: i128,
}

pub struct PhaseCarrierSpec {
    pub carrier_id: StableId,
    pub omega: RationalFixed,
    pub theta0: RationalFixed,
    pub quantization: PhaseQuantization,
    pub epoch_policy: EpochPolicy,
    pub role: PhaseCarrierRole,
}

pub enum PhaseCarrierRole {
    PrimaryClock,
    IntrinsicPhase,
    ExtrinsicPhase,
    SeamClock,
    HorizonClock,
    FinalizationClock,
    AddressShuffle,
}

pub enum PhaseQuantization {
    ContinuousFixed,
    KPolar { k: u32 },
    Tripolar { low: Fixed, high: Fixed },
    Quadrupolar,
}

pub enum EpochPolicy {
    FloorTurns,
    SaturatingTurns { max_epoch: u64 },
    Windowed { period: u64 },
}
\end{lstlisting}

\begin{lstlisting}
pub struct HypertorusSchedulerSpec {
    pub scheduler_id: StableId,
    pub carriers: Vec<PhaseCarrierSpec>,
    pub collision_policy: CollisionPolicy,
    pub address_policy: AddressPolicy,
    pub replay_policy: ReplayPolicy,
}

pub enum CollisionPolicy { Reject, StableTieBreakByHash, AllowWithDiagnostics }
pub enum AddressPolicy { Sequential, IrrationalPermutation, CarrierProduct }
pub enum ReplayPolicy { ByteIdenticalRequired, DiagnosticOnly }
\end{lstlisting}

\subsection{Scheduler-Zustand}
\begin{lstlisting}
pub struct PhaseCarrierState {
    pub carrier_id: StableId,
    pub step: u64,
    pub theta: Fixed,
    pub theta_mod_tau: Fixed,
    pub epoch: u64,
    pub quantized: PhaseBin,
}

pub enum PhaseBin {
    Continuous(Fixed),
    KPolar(u32),
    Tripolar(i8),
    Quadrupolar(i8, i8),
}

pub struct HypertorusSchedulerState {
    pub state_id: Hash256,
    pub spec_hash: Hash256,
    pub step: u64,
    pub carriers: Vec<PhaseCarrierState>,
    pub address_index: Option<u64>,
    pub diagnostics: Vec<SchedulerDiagnostic>,
}
\end{lstlisting}

\subsection{Normative Anforderungen}
\begin{enumerate}
  \item Carrier \MUST{} nach \code{carrier_id} sortiert canonicalisiert werden.
  \item Phasenfortschreibung \MUST{} deterministisch aus \code{step}, \code{omega}, \code{theta0} und \code{epoch_policy} berechnet werden.
  \item Scheduler-Entscheidungen \MUST{} replayfaehig sein.
  \item Nicht-wiederholende Shuffles \MUST{} deterministisch sein; echte Zufallsquellen sind im Audit-Pfad unzulaessig.
\end{enumerate}

\section{Phase Horizon und Event-Horizon-Gate}

\subsection{Datenstrukturen}
\begin{lstlisting}
pub struct PhaseHorizonWindowSpec {
    pub horizon_id: StableId,
    pub carrier_id: StableId,
    pub phase_start: Fixed,
    pub phase_end: Fixed,
    pub allowed_epochs: Option<Vec<u64>>,
    pub duty_cycle_min: Fixed,
    pub duty_cycle_max: Fixed,
    pub required_alignment: Fixed,
}

pub struct PhaseHorizonReport {
    pub report_id: Hash256,
    pub horizon_id: StableId,
    pub carrier_id: StableId,
    pub step: u64,
    pub epoch: u64,
    pub theta_mod_tau: Fixed,
    pub open: bool,
    pub alignment_score: Fixed,
    pub jitter_score: Fixed,
    pub passed: bool,
    pub evidence_refs: Vec<EvidenceRef>,
}
\end{lstlisting}

\subsection{Gate-Regel}
Ein Horizon-Fenster ist offen, wenn:
\[
G_{horizon}=1 \Longleftrightarrow
\widetilde{\theta}_i(k)\in [\alpha,\beta]
\wedge \operatorname{epoch}_i(k)\in E_{allowed}
\wedge A_h(k)\ge \theta_h
\wedge J_h(k)\le \epsilon_j.
\]

Falls \code{HorizonPolicy = StrictWindow}, \MUST{} jede Finalisierung bei geschlossenem Fenster scheitern und \code{Hold} erzeugen.

\section{Exclusion Axis}

\subsection{Bedeutung}
Die ExclusionAxis ist die neutrale Implementierung der Realitaetssicherung: instabile, nichtkohärente, nicht-rueckfuehrbare oder kontinutitaetsverletzende Pfade werden entfernt, bevor Expansion oder Finalisierung erlaubt ist.

\subsection{Datenstrukturen}
\begin{lstlisting}
pub struct ExclusionAxisReport {
    pub report_id: Hash256,
    pub por_score: Fixed,
    pub telescope_focus_score: Fixed,
    pub reality_cut_passed: bool,
    pub excision_count: u64,
    pub retained_count: u64,
    pub exclusion_score: Fixed,
    pub passed: bool,
    pub evidence_refs: Vec<EvidenceRef>,
}
\end{lstlisting}

\subsection{Gate-Regel}
\[
G_{exclusion}=1 \Longleftrightarrow
PoR=1 \wedge T_\gamma\text{ stabil} \wedge RealityCut=1 \wedge Score_E\ge \theta_E.
\]

\section{Expansion Axis}

\subsection{Bedeutung}
Die ExpansionAxis erzeugt aus einem gueltigen, stabilisierten Zustand ein gerichtetes Handlungspotential. Sie ersetzt keine Emission; sie berechnet nur, ob der Zustand ein finales Output-Ereignis tragen kann.

\subsection{Resonanzkanal}
Jeder Kanal \(i\) besitzt:
\[
D_i(k)=\psi_i(k)\cdot\rho_i(k)\cdot\omega_i(k)\cdot\Lambda_i(k).
\]
Hier sind \(\psi\), \(\rho\), \(\omega\) neutrale Scores fuer Bedeutung/Signifikanz, Strukturkohärenz und Phasenbereitschaft. \(\Lambda_i\) ist eine deterministische Huelle, z. B. aus einem PhaseCarrier.

\subsection{Gesamtdynamik}
\[
D_{total}(k)=\left(\prod_{i=1}^{N}D_i(k)\right)\cdot\Omega(k).
\]
Die Implementierung \MUST{} eine overflow-sichere, fixed-point-kompatible Aggregation verwenden. Direktes Produkt \MAY{} durch Log-Summe oder saturierende Multiplikation ersetzt werden, sofern das Verfahren dokumentiert und replayfaehig ist.

\subsection{Datenstrukturen}
\begin{lstlisting}
pub struct ResonanceTriple {
    pub semantic_density: Fixed,   // psi
    pub structural_coherence: Fixed, // rho
    pub phase_readiness: Fixed,    // omega
}

pub struct ResonanceChannel {
    pub channel_id: StableId,
    pub triple: ResonanceTriple,
    pub envelope: Fixed,
    pub channel_score: Fixed,
    pub evidence_refs: Vec<EvidenceRef>,
}

pub struct ExpansionAxisReport {
    pub report_id: Hash256,
    pub channels: Vec<ResonanceChannel>,
    pub total_dynamic_score: Fixed,
    pub threshold: Fixed,
    pub spike: bool,
    pub expansion_score: Fixed,
    pub passed: bool,
    pub evidence_refs: Vec<EvidenceRef>,
}
\end{lstlisting}

\subsection{Gate-Regel}
\[
G_{expansion}=1 \Longleftrightarrow D_{total}(k)>\Theta(k) \wedge Score_X\ge \theta_X.
\]

\section{Finalization Core}

\subsection{Bedeutung}
Der FinalizationCore fuehrt gueltige Projektion und Expansion in einen irreversiblen, content-addressed Handlungsknoten zusammen. Irreversibel bedeutet hier: Nach Erzeugung des finalen Artefakts darf dessen Hash nicht ohne neuen RunDescriptor veraendert werden.

\subsection{Konvergenzfeld}
Ein finales Konvergenzfeld wird als aggregierter Vektor modelliert:
\[
C(k)=\operatorname{canon\_aggregate}(V_{view},V_{seam},V_{witness},V_{quality},V_{horizon}).
\]

\subsection{Action Vector}
\begin{lstlisting}
pub struct ActionVector {
    pub vector_id: Hash256,
    pub source_bundle_id: Hash256,
    pub convergence_hash: Hash256,
    pub direction_hash: Hash256,
    pub magnitude: Fixed,
    pub target_kind: String,
    pub evidence_refs: Vec<EvidenceRef>,
}
\end{lstlisting}

\subsection{Finalization Report}
\begin{lstlisting}
pub struct FinalizationCoreReport {
    pub report_id: Hash256,
    pub exclusion_passed: bool,
    pub expansion_passed: bool,
    pub convergence_score: Fixed,
    pub finalization_score: Fixed,
    pub action_vector: Option<ActionVector>,
    pub passed: bool,
    pub evidence_refs: Vec<EvidenceRef>,
}
\end{lstlisting}

\subsection{Gate-Regel}
\[
G_{finalization}=1 \Longleftrightarrow
G_{exclusion}=1 \wedge G_{expansion}=1 \wedge K_C(k)\ge \theta_C \wedge Score_F\ge \theta_F.
\]

\section{Holistic Matrix Gate}

\subsection{Matrixlogik}
Die finale Matrix ist:
\[
M(k)=
\begin{cases}
E(k), & G_{exclusion}=1 \wedge G_{expansion}=1 \wedge G_{finalization}=1,\\
\varnothing, & \text{sonst.}
\end{cases}
\]

\subsection{Gesamtgate}
Die v0.2-Gesamtbedingung lautet:
\[
G_{HGF}=1 \Longleftrightarrow
G_{UPF}=1
\wedge G_{horizon}=1
\wedge G_{exclusion}=1
\wedge G_{expansion}=1
\wedge G_{finalization}=1
\wedge ReplayReady=1.
\]

\code{G_UPF} umfasst NullCenter, MirrorSeam, Telescope, Loopback, TopologicalWitness, ProjectionQuality, RealityCut und PathExcision-Sicherheit aus v0.1.

\subsection{Report}
\begin{lstlisting}
pub struct HolisticMatrixGateReport {
    pub report_id: Hash256,
    pub upf_passed: bool,
    pub horizon_passed: bool,
    pub exclusion_passed: bool,
    pub expansion_passed: bool,
    pub finalization_passed: bool,
    pub replay_ready: bool,
    pub passed: bool,
    pub failure_policy: FailurePolicyV2,
    pub reasons: Vec<String>,
    pub evidence_refs: Vec<EvidenceRef>,
}

pub enum FailurePolicyV2 {
    Hold,
    RefineProjection,
    SplitPhase,
    MigrateCarrier,
    WaitForHorizon,
    StrengthenExclusion,
    SuppressExpansion,
    RecomputeFinalization,
    RequireHumanReview,
    Abort,
}
\end{lstlisting}

\section{Finalized Emission}

\subsection{Datenstruktur}
\begin{lstlisting}
pub struct FinalizedEmission {
    pub emission_id: Hash256,
    pub rd_hash: Hash256,
    pub source_projection_bundle_id: Hash256,
    pub null_center_id: Hash256,
    pub scheduler_state_hash: Hash256,
    pub horizon_report_hash: Hash256,
    pub exclusion_report_hash: Hash256,
    pub expansion_report_hash: Hash256,
    pub finalization_report_hash: Hash256,
    pub holistic_gate_report_hash: Hash256,
    pub action_vector: ActionVector,
    pub target_artifacts: Vec<TargetArtifact>,
    pub ledger_entry_hash: Hash256,
    pub evidence_refs: Vec<EvidenceRef>,
}
\end{lstlisting}

\subsection{Normative Anforderungen}
\begin{enumerate}
  \item \code{FinalizedEmission} \MUST{} nur bei bestandenem \code{HolisticMatrixGate} erzeugt werden.
  \item \code{FinalizedEmission} \MUST{} alle Hashes enthalten, die zur Wiederholung der Entscheidung erforderlich sind.
  \item \code{FinalizedEmission} \MUSTNOT{} als \code{SemanticCrystal} serialisiert werden.
  \item \code{FinalizedEmission} \MAY{} an die PSE-Bridge uebergeben werden, falls Feature \code{projection-pse-bridge} aktiv ist.
\end{enumerate}

\section{Projektions-Pipeline v0.2}

\subsection{Referenzablauf}
\begin{enumerate}
  \item Lade \code{ProjectionRunDescriptorV2} und internen Zustand.
  \item Canonicalize Input und berechne \code{NullCenterObject}.
  \item Erzeuge \code{HypertorusSchedulerState} fuer Step \(k\).
  \item Evaluiere \code{PhaseHorizonWindow}.
  \item Fuehre UPF-v0.1-Kernpipeline aus: Mount, Telescope, Loopback, W/K/T/P, MirrorSeam, Witness, RealityCut, Quality, ProjectionGate.
  \item Erzeuge bei UPF-Gate-Pass ein \code{ProjectionEmissionBundle}; bei Fail ein \code{ProjectionHoldReport}.
  \item Evaluiere \code{ExclusionAxisReport}.
  \item Evaluiere \code{ExpansionAxisReport}.
  \item Evaluiere \code{FinalizationCoreReport}.
  \item Evaluiere \code{HolisticMatrixGateReport}.
  \item Bei Gate-Pass: erzeuge \code{FinalizedEmission} und \code{FinalizationCertificate}.
  \item Bei Gate-Fail: erzeuge \code{ProjectionHoldReport} mit v0.2-FailurePolicy.
  \item Schreibe Ledger, Manifest und Replay-Material.
\end{enumerate}

\subsection{Pseudocode}
\begin{lstlisting}
fn run_projection_v2(input: ProjectionInput, rd: ProjectionRunDescriptorV2)
    -> ProjectionV2Outcome
{
    let canonical = canonicalize(input, rd.canonicalization_version);
    let null_center = null_center_project(&canonical, &rd);

    let scheduler = hypertorus_step(&rd.hypertorus, rd.step());
    let horizon = evaluate_phase_horizon(&scheduler, &rd);

    let upf = run_upf_v1_core(canonical, null_center, &scheduler, &rd);
    if !upf.gate_passed() {
        return ProjectionV2Outcome::Hold(make_v2_hold(upf, horizon, &rd));
    }

    if !horizon.passed {
        return ProjectionV2Outcome::HoldWithPolicy(
            make_wait_for_horizon_hold(upf, horizon, &rd)
        );
    }

    let exclusion = evaluate_exclusion_axis(&upf, &rd);
    let expansion = evaluate_expansion_axis(&upf, &scheduler, &rd);
    let finalization = evaluate_finalization_core(&upf, &exclusion, &expansion, &rd);

    let matrix_gate = holistic_matrix_gate(
        &upf, &horizon, &exclusion, &expansion, &finalization, &rd
    );

    if matrix_gate.passed {
        ProjectionV2Outcome::Finalized(make_finalized_emission(
            upf, scheduler, horizon, exclusion, expansion, finalization, matrix_gate, rd
        ))
    } else {
        ProjectionV2Outcome::Hold(make_v2_hold_from_matrix(
            upf, horizon, exclusion, expansion, finalization, matrix_gate, rd
        ))
    }
}
\end{lstlisting}

\section{Reports und Zertifikate}

\subsection{Outcome}
\begin{lstlisting}
pub enum ProjectionV2Outcome {
    Finalized(FinalizedEmission),
    ProjectionOnly(ProjectionEmissionBundle),
    Hold(ProjectionHoldReportV2),
    WaitForHorizon(ProjectionHoldReportV2),
    ExcisionOnly(RealityCutReport),
    NeedsCarrierMigration(ProjectionHoldReportV2),
    InvalidInput(ProjectionDiagnostic),
    DeterminismViolation(ProjectionDiagnostic),
}
\end{lstlisting}

\subsection{HoldReport v0.2}
\begin{lstlisting}
pub struct ProjectionHoldReportV2 {
    pub hold_id: Hash256,
    pub rd_hash: Hash256,
    pub source_state_id: Hash256,
    pub scheduler_state_hash: Hash256,
    pub horizon_report: PhaseHorizonReport,
    pub upf_gate_report: ProjectionGateReport,
    pub matrix_gate_report: HolisticMatrixGateReport,
    pub diagnostics: Vec<ProjectionDiagnostic>,
    pub recommended_next_policy: Option<FailurePolicyV2>,
}
\end{lstlisting}

\subsection{FinalizationCertificate}
\begin{lstlisting}
pub struct FinalizationCertificate {
    pub certificate_id: Hash256,
    pub rd_hash: Hash256,
    pub input_hash: Hash256,
    pub null_center_id: Hash256,
    pub scheduler_state_hash: Hash256,
    pub horizon_report_hash: Hash256,
    pub projection_certificate_hash: Hash256,
    pub exclusion_report_hash: Hash256,
    pub expansion_report_hash: Hash256,
    pub finalization_report_hash: Hash256,
    pub holistic_gate_report_hash: Hash256,
    pub finalized_emission_hash: Hash256,
}
\end{lstlisting}

\section{CLI-Schnittstelle}

\subsection{Subcommands}
\begin{lstlisting}
pse-traverse-projection inspect <input.json>
pse-traverse-projection schedule <input.json> --rd rd.json --step <k> --out scheduler.json
pse-traverse-projection horizon <scheduler.json> --rd rd.json --out horizon.json
pse-traverse-projection project <input.json> --rd rd.json --out projection.json
pse-traverse-projection gate <projection.json> --out gate.json
pse-traverse-projection emit <projection.json> --out bundle.json
pse-traverse-projection finalize <bundle.json> --rd rd.json --out final.json
pse-traverse-projection replay <bundle-or-final-or-hold.json>
pse-traverse-projection verify <certificate.json>
\end{lstlisting}

\subsection{CLI-Anforderungen}
\begin{itemize}
  \item \code{schedule} \MUST{} SchedulerState byte-identisch aus RD und Step erzeugen.
  \item \code{horizon} \MUST{} PhaseHorizonReport erzeugen, ohne Projektion oder Emission auszufuehren.
  \item \code{project} \MUST{} die UPF-Kernpipeline bis ProjectionGateReport ausfuehren.
  \item \code{emit} \MUST{} fehlschlagen, wenn UPF-Gate nicht bestanden ist.
  \item \code{finalize} \MUST{} fehlschlagen, wenn Horizon, Exclusion, Expansion oder Finalization nicht bestanden ist.
  \item \code{replay} \MUST{} kanonische Byte-Identitaet pruefen.
  \item \code{verify} \MUST{} alle Zertifikatshashes, RD-Hash, SchedulerState, HorizonReport, MatrixGateReport und Emission validieren.
\end{itemize}

\section{Determinismus und Canonicalization}

\subsection{Audit-Pfad}
Der v0.2-Audit-Pfad umfasst:
\begin{itemize}
  \item \code{ProjectionRunDescriptorV2};
  \item canonical input bytes;
  \item \code{NullCenterObject};
  \item \code{HypertorusSchedulerSpec} und \code{HypertorusSchedulerState};
  \item \code{PhaseHorizonReport};
  \item ProjectionLayerSpec-Liste und Mount-Reihenfolge;
  \item TelescopeView und LoopbackReport;
  \item MirrorSeamReports;
  \item TopologicalWitness;
  \item RealityCutReport;
  \item ProjectionGateReport;
  \item ExclusionAxisReport;
  \item ExpansionAxisReport;
  \item FinalizationCoreReport;
  \item HolisticMatrixGateReport;
  \item FinalizedEmission oder HoldReport.
\end{itemize}

Im Audit-Pfad sind verboten:
\begin{itemize}
  \item Wall-clock timestamps;
  \item Plattform-Floats;
  \item ungeordnete HashMap-Iteration;
  \item OS-abhaengige Pfadnormalisierung;
  \item nicht gespeicherte RNG-Entscheidungen;
  \item externe Netzwerkzustaende;
  \item unstabile Approximationen irrationaler Frequenzen.
\end{itemize}

\section{Integration mit bestehenden Crates}

\subsection{Datei- und Modulzuschnitt}
\begin{lstlisting}
crates/pse-traverse/src/projection/
  mod.rs
  null_center.rs
  projection_layer.rs
  mount.rs
  telescope.rs
  loopback.rs
  mirror_seam.rs
  phase_carrier.rs
  phase_ladder.rs
  hypertorus_scheduler.rs
  phase_horizon.rs
  resonance_triple.rs
  exclusion_axis.rs
  expansion_axis.rs
  finalization_core.rs
  holistic_matrix_gate.rs
  finalized_emission.rs
  consensus_cycle.rs
  path_excision.rs
  nullpoint_detection.rs
  topology_witness.rs
  projection_quality.rs
  reality_cut.rs
  gated_emission.rs
  report.rs
  certificate.rs
  replay.rs

tools/pse-traverse-projection-cli/src/main.rs
\end{lstlisting}

\subsection{Feature Flags}
\begin{longtable}{p{0.28\linewidth}p{0.62\linewidth}}
\toprule
Feature & Bedeutung \\
\midrule
\endhead
\code{projection} & Aktiviert Projektionskern. \\
\code{projection-cli} & Baut CLI-Tool. \\
\code{projection-topology} & Aktiviert TopologicalWitness-Schicht. \\
\code{projection-hypertorus} & Aktiviert HypertorusScheduler, PhaseCarrier und HorizonWindow. \\
\code{projection-finalization} & Aktiviert Exclusion/Expansion/Finalization-Matrix und FinalizedEmission. \\
\code{projection-pse-bridge} & Erlaubt Weitergabe finalisierter Emissionen an bestehende PSE-Bridge. \\
\code{projection-adapters} & Aktiviert optionale DomainAdapter; Kern bleibt domänenneutral. \\
\bottomrule
\end{longtable}

\section{Fehlersemantik}

\subsection{Fail-Closed-Regel}
Jede Unsicherheit in NullCenter, Scheduler, Horizon, Seam, Witness, RealityCut, Expansion, Finalization, Canonicalization oder Replay \MUST{} zu \code{Hold} oder haerterem Fail fuehren. Kein unsicheres Ergebnis darf emittiert oder finalisiert werden.

\subsection{Migration bei Phase-Failure}
Bei Horizon-Failure \SHOULD{} das System nicht blind abbrechen. Es \SHOULD{} eine maschinenlesbare Empfehlung erzeugen:
\begin{itemize}
  \item \code{WaitForHorizon}, falls der naechste gueltige Step berechenbar ist;
  \item \code{MigrateCarrier}, falls alternative Carrier im RD deklariert sind;
  \item \code{SplitPhase}, falls Jitter oder Kollisionsdiagnostik ein Phasenproblem zeigt;
  \item \code{Abort}, falls RD oder Scheduler inkonsistent ist.
\end{itemize}

\section{Testplan}

\subsection{Unit Tests}
\begin{enumerate}
  \item \code{hypertorus_scheduler_is_deterministic}
  \item \code{phase_carriers_are_sorted_before_hashing}
  \item \code{phase_epoch_counter_matches_turns}
  \item \code{phase_window_opens_only_in_declared_interval}
  \item \code{phase_window_rejects_wrong_epoch}
  \item \code{k_polar_quantization_is_stable}
  \item \code{tripolar_dead_zone_is_stable}
  \item \code{address_shuffle_replays_identically}
  \item \code{horizon_gate_fails_closed}
  \item \code{exclusion_without_por_fails}
  \item \code{expansion_without_spike_fails}
  \item \code{finalization_without_convergence_fails}
  \item \code{holistic_matrix_requires_all_axes}
  \item \code{finalized_emission_only_after_matrix_gate}
  \item \code{finalized_emission_changes_when_epoch_changes}
  \item \code{finalization_certificate_changes_when_action_vector_changes}
\end{enumerate}

\subsection{Integration Tests}
\begin{enumerate}
  \item \code{cli_schedule_writes_scheduler_state}
  \item \code{cli_horizon_writes_horizon_report}
  \item \code{cli_finalize_refuses_closed_horizon}
  \item \code{cli_finalize_refuses_failed_exclusion}
  \item \code{cli_finalize_refuses_failed_expansion}
  \item \code{cli_finalize_refuses_failed_finalization}
  \item \code{cli_finalize_writes_finalized_emission}
  \item \code{cli_replay_byte_identical_finalized_emission}
  \item \code{cli_verify_detects_modified_scheduler_state}
  \item \code{pse_bridge_receives_finalized_emission_but_not_crystal}
\end{enumerate}

\subsection{Property Tests}
\begin{enumerate}
  \item Permutierte Carrier-Eingaben erzeugen identische SchedulerSpec-Hashes.
  \item Gate-Pass ist monoton falsch bei Verschaerfung einzelner Schwellen, sofern Messwerte unveraendert bleiben.
  \item HorizonReports sind Funktionen von RD, Step und SchedulerSpec.
  \item FinalizedEmission-ID ist Funktion aller referenzierten Hashes.
  \item Geschlossene HorizonWindows koennen keine FinalizedEmission erzeugen.
  \item Ein einzelnes fehlgeschlagenes Pflichtgate verhindert Finalisierung.
\end{enumerate}

\section{Definition of Done}
Eine Implementierung von \code{PSE-TRAVERSE-PROJECTION-02} ist abgeschlossen, wenn:
\begin{enumerate}
  \item \code{cargo test --workspace} besteht.
  \item Der Projektionskern ohne Netzwerkabhaengigkeit kompiliert.
  \item Alle oeffentlichen Typen dokumentiert sind.
  \item Alle Gates fail-closed sind.
  \item Kein gate-relevanter Wert Plattform-Floats im Audit-Pfad nutzt.
  \item \code{NullCenterObject} content-addressed und replayfaehig ist.
  \item \code{HypertorusSchedulerState} byte-identisch replayfaehig ist.
  \item \code{PhaseHorizonWindow} Finalisierung bei geschlossenem Fenster verhindert.
  \item \code{ExclusionAxisReport}, \code{ExpansionAxisReport} und \code{FinalizationCoreReport} gate-wirksam sind.
  \item \code{HolisticMatrixGate} bei einem einzigen fehlgeschlagenen Pflichtgate Hold erzeugt.
  \item \code{FinalizedEmission} nur bei Gate-Pass erzeugt wird.
  \item \code{ProjectionHoldReportV2} maschinenlesbare Gruende und FailurePolicy enthaelt.
  \item \code{RealityCut} instabile Pfade deterministisch exzidiert.
  \item \code{TopologicalWitness} optional konfigurierbar, bei Aktivierung aber gate-wirksam ist.
  \item \code{PathExcisionJump} nur bei Kairos-/Gate-Pass erlaubt ist.
  \item CLI \code{schedule}, \code{horizon}, \code{project}, \code{emit}, \code{finalize}, \code{replay}, \code{verify} end-to-end funktioniert.
  \item Zwei identische Runs byte-identische Reports, Bundles und FinalizedEmission-IDs erzeugen.
  \item Certificate-Verification veraenderte RD, SchedulerState, Reports oder Artifacts erkennt.
  \item PSE-Bridge einziger Pfad zu \code{SemanticCrystal} bleibt.
\end{enumerate}

\section{Coding-Agent-Auftrag}

\begin{lstlisting}
PATCH-ID: PSE-TRAVERSE-PROJECTION-02
TITLE: Hypertorus-Gated Finalization Layer

GOAL:
Implementiere die v0.2-Gesamtschicht fuer pse-traverse-projection.
Die Schicht ersetzt v0.1 normativ, indem sie alle v0.1-Funktionen
abdeckt und um HypertorusScheduler, PhaseHorizonWindow,
ExclusionAxis, ExpansionAxis, FinalizationCore, HolisticMatrixGate
und FinalizedEmission erweitert.

MUST:
- neue Module unter crates/pse-traverse/src/projection/ anlegen;
- ProjectionRunDescriptorV2, ProjectionThresholdsV2 und Policies implementieren;
- HypertorusSchedulerSpec/State deterministisch und replayfaehig implementieren;
- PhaseHorizonWindow und HorizonReport implementieren;
- ExclusionAxisReport, ExpansionAxisReport und FinalizationCoreReport implementieren;
- HolisticMatrixGate fail-closed implementieren;
- FinalizedEmission und FinalizationCertificate content-addressed erzeugen;
- CLI um schedule, horizon und finalize erweitern;
- Tests gemaess Testplan liefern;
- PSE SemanticCrystal-Erzeugung ausschliesslich bestehender PSE-Bridge ueberlassen.

MUST NOT:
- Kosmologie, Zeta, Riemann, Trading, Quantum-Hardware oder Hardware-Fahrzeugsemantik in den Core einbauen;
- nichtdeterministische RNG- oder Wall-clock-Werte in Digests schreiben;
- Plattform-Floats fuer Gate-/Audit-Hashes verwenden;
- bei fehlgeschlagenem Horizon/Gate finalisieren;
- SemanticCrystal in dieser Schicht erzeugen.

SHOULD:
- BTreeMap/BTreeSet fuer kanonische Strukturen verwenden;
- Feature Flags projection-hypertorus und projection-finalization anlegen;
- goldene Replay-Fixtures fuer offene/geschlossene HorizonWindows liefern;
- v0.1-Kompatibilitaet ueber Adapter oder Re-Exports abdecken.
\end{lstlisting}

\section{Referenz-Zustandsmaschine}

\begin{longtable}{p{0.22\linewidth}p{0.23\linewidth}p{0.45\linewidth}}
\toprule
Von & Nach & Bedingung \\
\midrule
\endhead
Idle & Loading & Input und RD vorhanden. \\
Loading & Canonicalizing & Dateien gelesen und Schema gueltig. \\
Canonicalizing & NullProjecting & Canonical bytes erzeugt. \\
NullProjecting & Scheduling & NullCenterObject erzeugt. \\
Scheduling & HorizonEvaluating & SchedulerState erzeugt. \\
HorizonEvaluating & Mounting & Horizon offen oder Policy erlaubt ProjectionOnly. \\
HorizonEvaluating & Holding & Horizon geschlossen und Finalisierung verlangt. \\
Mounting & Focusing & Alle ProjectionLayers gemountet. \\
Focusing & LoopbackTesting & TelescopeView bestimmt. \\
LoopbackTesting & Condensing & Loopback instabil. \\
LoopbackTesting & SeamEvaluating & Loopback stabil. \\
Condensing & SeamEvaluating & W/K/T/P-Zyklus abgeschlossen. \\
SeamEvaluating & Witnessing & Alle SeamReports erzeugt. \\
Witnessing & RealityCutting & Witness optional oder bestanden. \\
RealityCutting & QualityScoring & Pfade exzidiert/erhalten. \\
QualityScoring & ProjectionGating & QualityCascade berechnet. \\
ProjectionGating & ExclusionEvaluating & UPF-Gate bestanden. \\
ProjectionGating & Holding & UPF-Gate fehlgeschlagen. \\
ExclusionEvaluating & ExpansionEvaluating & Exclusion bestanden. \\
ExclusionEvaluating & Holding & Exclusion fehlgeschlagen. \\
ExpansionEvaluating & FinalizationEvaluating & Expansion bestanden. \\
ExpansionEvaluating & Holding & Expansion fehlgeschlagen. \\
FinalizationEvaluating & MatrixGating & Finalization bestanden. \\
FinalizationEvaluating & Holding & Finalization fehlgeschlagen. \\
MatrixGating & Finalized & HolisticMatrixGate bestanden. \\
MatrixGating & Holding & MatrixGate fehlgeschlagen. \\
Finalized & Certified & FinalizedEmission und Certificate erzeugt. \\
Holding & Certified & HoldReport und Certificate erzeugt. \\
Certified & Replayed & Replay erfolgreich. \\
jede Phase & Failed & Schema-, Determinismus-, Invarianz- oder Hashverletzung. \\
\bottomrule
\end{longtable}

\section{Schlussformel}
Die v0.2-Bedingung lautet:
\[
HGF(s,RD,k)=FinalizedEmission(s)
\]
\[
\text{nur wenn}\quad
\delta(s)\text{ kanonisch}
\wedge G_{UPF}=1
\wedge G_{horizon}=1
\wedge G_{exclusion}=1
\wedge G_{expansion}=1
\wedge G_{finalization}=1
\wedge ReplayReady=1.
\]

Andernfalls gilt:
\[
HGF(s,RD,k)=Hold(s,diagnostics).
\]

\end{document}
